\documentclass{report}
\usepackage{cmt}

\title{\huge{Algorithmische Mathematik I \\ Übungszettel 2}}
\author{\huge{Amy, Rojan}}
\date{2026-10-22}

\begin{document}

\maketitle
\tableofcontents

\chapter{Grundlagen}

\nt{
  Den Übungszettel gibt es
  \href{run:./pset_2026-10-22.pdf}{hier}.
}

\section{Landau-Symbole}

\dfn{}{
  Sei \(g\colon\bbN\to\bbR_{\ge 0}\). Dann definieren wir:
  \begin{align*}
    O(g)
    &:=
    \left\{
      f\colon\bbN\to\bbR_{\ge 0}
      \;\middle|\;
      \exists \alpha>0\,\exists n_0\in\bbN\,
      \forall n\ge n_0:
      f(n)\le \alpha g(n)
    \right\},\\
    \Omega(g)
    &:=
    \left\{
      f\colon\bbN\to\bbR_{\ge 0}
      \;\middle|\;
      \exists \alpha>0\,\exists n_0\in\bbN\,
      \forall n\ge n_0:
      f(n)\ge \alpha g(n)
    \right\},\\
    \Theta(g)&:=O(g)\cap\Omega(g).
  \end{align*}
}

\section{Aufgaben}

\qs{}{
  Beweisen Sie folgende Aussagen:
  \begin{enumerate}
    \item Für alle Funktionen
      \(f,g\colon\bbN\to\bbR_{\ge 0}\) gilt:
      \[
        f\in O(g)\iff g\in\Omega(f).
      \]
    \item Für alle Funktionen
      \(f,g,h\colon\bbN\to\bbR_{\ge 0}\) gilt:
      \[
        f\in O(g)\land g\in O(h)
        \implies f\in O(h).
      \]
  \end{enumerate}
}

\begin{myproof}
  \textbf{1. Hinrichtung:}
  Sei \(f\in O(g)\). Dann existieren \(\alpha>0\) und
  \(n_0\in\bbN\), sodass für alle \(n\ge n_0\) gilt:
  \[
    f(n)\le \alpha g(n).
  \]
  Durch Multiplikation mit \(1/\alpha>0\) folgt
  \[
    g(n)\ge \frac{1}{\alpha}f(n).
  \]
  Also gilt \(g\in\Omega(f)\).

  \textbf{Rückrichtung:}
  Sei \(g\in\Omega(f)\). Dann existieren \(\beta>0\) und
  \(n_0\in\bbN\), sodass für alle \(n\ge n_0\) gilt:
  \[
    g(n)\ge \beta f(n).
  \]
  Durch Division durch \(\beta>0\) erhalten wir
  \[
    f(n)\le \frac{1}{\beta}g(n).
  \]
  Also gilt \(f\in O(g)\).

  Damit ist die Äquivalenz bewiesen.

  \medskip
  \textbf{2.}
  Seien \(f\in O(g)\) und \(g\in O(h)\). Dann existieren
  \(\alpha,\beta>0\) und \(n_1,n_2\in\bbN\), sodass für alle
  \(n\ge n_1\) und \(n\ge n_2\) gilt:
  \[
    f(n)\le \alpha g(n)
    \qquad\text{und}\qquad
    g(n)\le \beta h(n).
  \]
  Setze \(n_0:=\max\{n_1,n_2\}\). Für alle \(n\ge n_0\) folgt
  \[
    f(n)\le \alpha g(n)
    \le \alpha\beta h(n).
  \]
  Da \(\alpha\beta>0\), gilt \(f\in O(h)\).
\end{myproof}

\qs{}{
  Beweisen Sie:
  \begin{enumerate}
    \item \(\log n=\Theta(\ln n)\)
    \item \(\log(n!)=\Theta(n\log n)\)
  \end{enumerate}
}

\begin{myproof}
  \textbf{1.}
  Da \(\log\) den Logarithmus zur Basis \(10\) bezeichnet, gilt
  \[
    \log n=\frac{1}{\ln 10}\ln n.
  \]
  Mit \(\alpha=1/\ln 10>0\) folgt unmittelbar
  \[
    \log n\in O(\ln n)
    \qquad\text{und}\qquad
    \log n\in\Omega(\ln n).
  \]
  Also gilt
  \[
    \log n\in\Theta(\ln n).
  \]

  \medskip
  \textbf{2.}
  Für alle \(n\ge 2\) gilt
  \[
    n!\le n^n.
  \]
  Wegen der Monotonie des Logarithmus folgt
  \[
    \log(n!)\le n\log n.
  \]
  Somit ist \(\log(n!)\in O(n\log n)\).

  Für die untere Schranke betrachten wir die letzten Faktoren
  des Produkts \(n!\). Für \(n\ge 4\) sind mindestens \(n/2\)
  Faktoren mindestens \(n/2\). Daher gilt
  \[
    n!\ge \left(\frac n2\right)^{n/2}.
  \]
  Durch Logarithmieren erhalten wir
  \[
    \log(n!)
    \ge \frac n2\log\left(\frac n2\right).
  \]
  Für \(n\ge 4\) gilt außerdem
  \[
    \log\left(\frac n2\right)\ge \frac12\log n.
  \]
  Folglich
  \[
    \log(n!)\ge \frac n4\log n.
  \]
  Damit ist \(\log(n!)\in\Omega(n\log n)\), und insgesamt gilt
  \[
    \log(n!)\in\Theta(n\log n).
  \]
\end{myproof}

\qs{}{
  Zeigen Sie, dass für alle \(k,l\in\bbN\) mit \(l\ge 1\) gilt:
  \begin{enumerate}
    \item \((\log n)^k=O(n^{1/l})\)
    \item \(2^{\sqrt{\log n}}=O(n^{1/l})\)
  \end{enumerate}
}

\begin{myproof}
  \textbf{1.}
  Setze \(x=\log n\). Da \(n=10^x\), genügt es zu zeigen, dass
  \[
    \frac{x^k}{10^{x/l}}\longrightarrow 0
    \qquad\text{für }x\to\infty.
  \]
  Die Exponentialfunktion wächst schneller als jedes Polynom.
  Daher konvergiert dieser Quotient gegen \(0\). Insbesondere
  existieren \(\alpha>0\) und \(n_0\in\bbN\), sodass für alle
  \(n\ge n_0\) gilt:
  \[
    (\log n)^k\le \alpha n^{1/l}.
  \]
  Also ist \((\log n)^k\in O(n^{1/l})\).

  \medskip
  \textbf{2.}
  Setze wieder \(x=\log n\). Für hinreichend große \(x\) gilt
  \[
    \sqrt{x}\log 2\le \frac{x\log 10}{l},
  \]
  denn \(x\) wächst schneller als \(\sqrt{x}\). Durch
  Exponenzieren folgt
  \[
    2^{\sqrt{x}}\le 10^{x/l}.
  \]
  Wegen \(x=\log n\) ist dies gleichbedeutend mit
  \[
    2^{\sqrt{\log n}}\le n^{1/l}
  \]
  für alle hinreichend großen \(n\). Also gilt
  \[
    2^{\sqrt{\log n}}\in O(n^{1/l}).
  \]
\end{myproof}

\qs{}{
  Zeigen Sie, dass die Potenzmenge von \(\bbN\), also die Menge
  aller Teilmengen von \(\bbN\), überabzählbar ist.
}

\begin{myproof}
  Wir führen einen Widerspruchsbeweis.

  Angenommen, \(\mathcal P(\bbN)\) sei abzählbar. Dann lassen
  sich alle Teilmengen von \(\bbN\) in einer Folge auflisten:
  \[
    \mathcal P(\bbN)=\{A_1,A_2,A_3,\ldots\}.
  \]
  Wir definieren die Menge
  \[
    B:=\{n\in\bbN\mid n\notin A_n\}.
  \]
  Da \(B\subseteq\bbN\) gilt, muss \(B=A_k\) für ein
  \(k\in\bbN\) sein.

  Wir unterscheiden zwei Fälle:
  \begin{itemize}
    \item Falls \(k\in A_k\), folgt aus der Definition von \(B\),
      dass \(k\notin B\). Wegen \(B=A_k\) ergibt sich
      \(k\notin A_k\), ein Widerspruch.
    \item Falls \(k\notin A_k\), folgt aus der Definition von \(B\),
      dass \(k\in B\). Wegen \(B=A_k\) ergibt sich
      \(k\in A_k\), ebenfalls ein Widerspruch.
  \end{itemize}
  In beiden Fällen entsteht ein Widerspruch. Daher ist
  \(\mathcal P(\bbN)\) überabzählbar.
\end{myproof}

\section{Programmierteil}
\qs{} {
Schreiben Sie ein C++-Programm, das alle Primzahlzwillinge (also Paare von Primzahlen, deren Differenz 2 ist) auflistet, die kleiner als 106 sind. Wie viele gibt es?
}

\codefile{naive.cpp}{C++}

\end{document}
